\begin{table}[!h]
\centering
\caption{\label{tab:regression_results}MICE-Pooled OLS Regression Results (Rubin's Rules, $M = 300$: 100 Python, 100 R, 100 Julia). Dep.\ var.: $\log(\text{Opportunity\_Cost})$.}
\centering
\resizebox{\ifdim\width>\linewidth\linewidth\else\width\fi}{!}{
\fontsize{9}{11}\selectfont
\begin{threeparttable}
\begin{tabular}[t]{lllllllllllll}
\toprule
\multicolumn{1}{c}{ } & \multicolumn{4}{c}{Python} & \multicolumn{4}{c}{R} & \multicolumn{4}{c}{Julia} \\
\cmidrule(l{3pt}r{3pt}){2-5} \cmidrule(l{3pt}r{3pt}){6-9} \cmidrule(l{3pt}r{3pt}){10-13}
Parameter & Coef. & SE & $p$ & Sig. & Coef. & SE & $p$ & Sig. & Coef. & SE & $p$ & Sig.\\
\midrule
Intercept & 12.283 & 0.038 & $<$ 0.001 & *** & 12.226 & 0.042 & $<$ 0.001 & *** & 12.247 & 0.039 & $<$ 0.001 & ***\\
Holes & 0.047 & 0.002 & $<$ 0.001 & *** & 0.053 & 0.003 & $<$ 0.001 & *** & 0.048 & 0.002 & $<$ 0.001 & ***\\
\bottomrule
\end{tabular}
\begin{tablenotes}
\item \textit{Note: } 
\item Sig. codes: *** $p < 0.001$. Rubin's Rules applied independently per language group ($M = 100$ each). Grand Mean = arithmetic mean of three independently pooled estimates.
\end{tablenotes}
\end{threeparttable}}
\end{table}
